% source: RKF/theorum/01_positive_eta_path_independence.md
\hypertarget{theorum-01-positive-eta-path-independence-by-amplitude-lift}{%
\chapter{Theorum 01: Positive-Eta Path Independence by Amplitude Lift}\label{theorum-01-positive-eta-path-independence-by-amplitude-lift}}

\hypertarget{source-provenance}{%
\section{Source provenance}\label{source-provenance}}

\begin{verbatim}
source repository: Parveen117/MP
primary PR:        #246 Reset RH route onto rotation-generated recognition topology
source branch:     agent/rh-rotation-metric-topology-reset
\end{verbatim}

\hypertarget{statement}{%
\section{Statement}\label{statement}}

On the fixed five-label event-response carrier, every positive-eta normalization/continuation rectangle is an identity loop.

Equivalently, for any admissible positive eta values \texttt{eta\_1,\ eta\_2},

\begin{verbatim}
normalize then eta-continue
=
eta-continue then normalize.
\end{verbatim}

No rank-one seam compensation is required for the positive-eta family itself.

\hypertarget{native-factorization}{%
\section{Native factorization}\label{native-factorization}}

The state-locked construction has the factorization

\begin{verbatim}
A_chart  = A U,
Theta_eta = Q_chart G_eta^(1/2),
G_eta = 1 + (A_+ + eta I_4),    eta > 0.
\end{verbatim}

Here:

\begin{verbatim}
A          raw native event amplitude, eta-independent;
U          one fixed unitary accepted chart;
Q_chart    the polar orientation of the charted amplitude;
G_eta      positive metric/scale factor for eta > 0.
\end{verbatim}

Because right multiplication by a unitary is polar-covariant and right multiplication by a strictly positive factor does not change the polar orientation,

\begin{verbatim}
polar(A U) = polar(A) U,
polar(Q_chart G_eta^(1/2)) = Q_chart.
\end{verbatim}

Thus the frame transport data are exactly

\begin{verbatim}
raw eta transport                 I_5
normalized eta transport          I_5
raw-to-normalized transport       U*
based loop                        I_5
\end{verbatim}

for every positive eta in the fixed five-label family.

\hypertarget{finite-state-locked-rectangle-evidence}{%
\section{Finite state-locked rectangle evidence}\label{finite-state-locked-rectangle-evidence}}

The user-local rectangle at

\begin{verbatim}
eta     = 0.010
eta + h = 0.011
\end{verbatim}

used the four actual amplitudes

\begin{verbatim}
raw(eta)          normalized(eta)
raw(eta+h)        normalized(eta+h)
\end{verbatim}

and reported

\begin{verbatim}
status                          PASS_STATE_LOCKED_UGD_ETA_NORMALIZATION_RECTANGLE_EVALUATION
loop classification             FINITE_ACTUAL_TWO_PATH_EQUALITY_PROVED
path transport difference       3.0669860554696173e-15
loop identity residual          3.1396725820594018e-15
defect phase                    1.0705180635847342e-40
defect-to-accepted leakage      1.7140530928543942e-17
accepted-block residual         3.139637946219698e-15
endpoint scale difference       0
maximum pairing return residual 6.071532165918825e-18
false checks                    none
\end{verbatim}

These numbers are the floating-point shadow of the exact positive-eta family theorem.

\hypertarget{what-this-closes}{%
\section{What this closes}\label{what-this-closes}}

\begin{verbatim}
positive-eta raw amplitude path equality       CLOSED
positive-eta normalization path equality       CLOSED
fixed five-label UGD rectangle                 CLOSED
need for new observer label at positive eta    REMOVED
rank-one seam compensation at positive eta     NOT NEEDED
\end{verbatim}

\hypertarget{what-this-does-not-close}{%
\section{What this does not close}\label{what-this-does-not-close}}

At eta zero,

\begin{verbatim}
G_0 = 1 + A_+.
\end{verbatim}

If \texttt{A\_+} has a kernel, the normalized response can lose support. Therefore the next eta problem is support/Fredholm descent, not another path-holonomy problem.

\begin{verbatim}
eta-zero support descent                  OPEN
kernel-line target visibility             OPEN
continuum UGD/Fredholm identification      OPEN
actual RH target map                       OPEN
odd/even or single-sector terminal sign    NOT DECIDED BY THIS THEOREM
\end{verbatim}

\hypertarget{use-in-the-framework}{%
\section{Use in the framework}\label{use-in-the-framework}}

Consume this theorem only as the positive-eta path-equivalence adapter:

\begin{verbatim}
positive eta family
-> no normalization/path-order ambiguity
-> fixed five-label carrier is lawful
-> eta-zero problem becomes support descent
\end{verbatim}

\begin{flushright}\small\texttt{RKF/theorum/01\_positive\_eta\_path\_independence.md}\end{flushright}
